\section{Ordinary-line tags and the final loss estimate}\label{sec:closing}

The ordinary-line tag \(r_\ell=\rho\) must be controlled uniformly
over all directions. We use fixed moment orders throughout this
section. Retaining the line's height and counting only the spatial
cells that it meets avoids a growing moment order or constants
uniform over all orders.

\begin{lemma}[Linewise cell count]\label{lem:cells}
A nonaxis line meets at most \(2G-1\) open cells of a fixed
\(G\times G\) grid. Its lattice points have at most two parity types.
Consequently there are at most \(C_{\rm line}=4G\) nonempty cell/type
pieces on it, uniformly in \(n\), root and direction.
\end{lemma}
\begin{proof}
Inside the open square a nonaxis line crosses each of the
\(2(G-1)\) internal grid lines at most once. These crossings divide
it into at most \(2G-1\) intervals. Convexity prevents re-entry
into a cell. Coincident crossings at a vertex only reduce this count.
The integer points are \(p+kv\), with types alternating between
\(r\) and \(r+\bar v\). Thus there are at most
\(2(2G-1)\le4G\) pieces. Their index set is determined before the
random process is run.
\end{proof}

For an ordinary line \(\ell\ni p\) and one such piece \(E\), put
\[
 H_{\ell,E}(t)=r_\ell(t)
       \sum_{\substack{q\in(\ell\cap E)\setminus\{p\}\\q\text{ active}}}w_q(t),
 \qquad L_{\ell,E}=| (\ell\cap E)\setminus\{p\}|.
\]
Stop the fixed-label hierarchy at the first arrival at
\(H_{\ell,E}\ge L_{\ell,E}+en\), or the root boundary, whichever
comes first. As usual, observations at a boundary precede changes
of positive-retirement flags. For fixed \(m\), let
\(\Theta_m(n)=(1-\epsilon_n)^{-\binom{m+1}2}\).
Use the common-stop form \eqref{eq:mixedstop} with deterministic
horizon \(T\). Expansion by the number \(a\) of distinct leaves gives
\begin{equation}\label{eq:tagmoment}
 \E^{(p)}H_{\ell,E}^m
 \le\Theta_m(n)\sum_{a=1}^m
 \left\{\begin{matrix}m\\a\end{matrix}\right\}
 B^{m-a}L_{\ell,E}^a.
\end{equation}
The braces are Stirling numbers of the second kind. For distinct
leaves the exact connection is \(r^a\); for repetitions use
\(r^m\le r^a\), \(w^j\le B^{j-1}w\). There are no witnesses and
no designated-line connections, so no extra fixed multiplier is
needed. Ordered distinct-leaf tuples are counted by
\((L)_a\le L^a\).

\begin{lemma}[Finite height sum]\label{lem:height}
For \(e>0\), integers \(m\ge3\) and \(N\ge0\),
\begin{equation}\label{eq:height}
 \sum_{h=1}^N4h\left(\frac2{2+eh}\right)^m
 \le\frac{16}{e^2}\left(\frac2{2+e}\right)^{m-3}.
\end{equation}
\end{lemma}
\begin{proof}
Put \(a=e/2\). For \(h\ge1\),
\[
 h(1+ah)^{-m}\le\frac{(1+a)^{-(m-3)}}{a(1+ah)^2},
\]
and
\[
 (1+ah)^{-2}\le\frac1a
 \left[\frac1{1+a(h-1)}-\frac1{1+ah}\right].
\]
Sum the latter telescoping difference. Multiplication by four
gives \eqref{eq:height}, since \(4/a^2=16/e^2\).
\end{proof}

Apply Markov to \eqref{eq:tagmoment}, sum first over the at most
\(C_{\rm line}\) pieces of each line, then over directions of
height \(h\), using \(4h\) directions and \(L\le2n/h\).
The all-distinct term is bounded by Lemma~\ref{lem:height}.
Bounding the denominator in each repeated term by \((en)^m\)
gives the explicit remainder
\begin{align}\label{eq:tagremainder}
 R_m(n;e)=4e^{-m}\bigg[&2\left\{\begin{matrix}m\\1\end{matrix}\right\}
                 n(B/n)^{m-1}
 +4\left\{\begin{matrix}m\\2\end{matrix}\right\}
                 (1+\log n)(B/n)^{m-2}\notag\\
 &+\sum_{a=3}^{m-1}2^a
 \left\{\begin{matrix}m\\a\end{matrix}\right\}
 \left(1+\frac1{a-2}\right)(B/n)^{m-a}\bigg].
\end{align}
Thus the probability of any such first crossing, under the fixed
root's point law, is at most
\begin{equation}\label{eq:tagprob}
 C_{\rm line}\Theta_m(n)
 \left[\frac{16}{e^2}\left(\frac2{2+e}\right)^{m-3}
       +R_m(n;e)\right].
\end{equation}
No union over times, or over all roots, has been used.
For sufficiently large fixed \(m\), \(R_m(n;e)\to0\); it is enough
that \(1-(1-d)(m-1)<0\) and \(m>2\).

On the complementary good event, sum the cell bounds after multiplying
by their deterministic mask values, then integrate \eqref{eq:tagflow}.
The self point costs a further \(TB/n\). The result is
\begin{equation}\label{eq:tagcapacity}
 1-r_\ell\le\text{deterministic ordinary-line capacity}
       +C_{\rm line}Te+TB/n+O(n^{-1}).
\end{equation}
Choose \(e=(1-\rho-\lambda)/(4TC_{\rm line})\) when \(T>0\).
The zero-time case has no births and is trivial. For large \(n\),
the last two terms are at most \((1-\rho-\lambda)/4\).
Using capacity at most \(\lambda\), the right side is then strictly
less than \(1-\rho\), precluding a tag. For any requested positive
error, choose \(m\) large enough that the first term of
\eqref{eq:tagprob} is below half that error, and then \(n\) large
enough for \(\Theta_m(n)\) and the remainder. This proves
\begin{equation}\label{eq:tagtail}
 \sup_p\Prob^{(p)}(\text{tag is an early saving cause})=o(1).
\end{equation}
This linewise count also applies to the designated-line clock
argument (an axis line has at most \(2G\) cell/type pieces).
It does not replace the full number of cells in the uniform
bounded-variation estimates of Section~\ref{sec:geometry}.

\begin{proof}[Completion of the analytic reduction]
The weight-cap estimate \eqref{eq:captail}, tag estimate
\eqref{eq:tagtail}, and global-stop estimate \eqref{eq:globaltail}
give \(\sup_p\Prob^{(p)}(E_p)=o(1)\). Their events may overlap;
a union bound suffices. Proposition~\ref{prop:loss} implies
\(\E|S_n|\ge(1-o(1))nA_n\), whereas the exact nonnegative loss
identity also gives \(\E|S_n|\le nA_n\). By Lemma~\ref{lem:chord},
\(A_n=A+O(n^{-1})\), proving Theorem~\ref{thm:bridge}.

For fixed \(c<A\) the expectation exceeds \(cn\) for all large
\(n\). If every possible legal output had cardinality smaller than
\(cn\), its expectation would be smaller too, a contradiction.
Equivalently, with \(k=\lceil cn\rceil\) and \(|S_n|\le2n\),
\[
 \Prob(|S_n|\ge k)\ge
 \frac{\E|S_n|-(k-1)}{2n-(k-1)}>0.
\]
This is an existence argument; it does not supply a practical
algorithm producing the target density on moderate grids.
\end{proof}

